Neural computer performance hinges on trace fidelity. We collect frames, text, and actions with precise timestamps so the model never has to infer when an action occurred. Commodity tooling keeps collection standardized and auditable. Early pilots showed that seemingly minor choices, such as fonts, sampling rate, and latency alignment, can shift OCR accuracy by up to 6~pp and pointer MAE by 4~px, so we lock these settings before training.

\noindent\textbf{CLI traces.}
For \nccligen{}, we prepare two CLI data modes. The first, CLIGen (General), starts from publicly released \texttt{asciinema}\footnote{\url{https://asciinema.org/}} \texttt{.cast} trajectories. We replay and render each session with the official \texttt{asciinema} library so palette transitions, cursor state, and terminal geometry are preserved; the underlying buffer provides text supervision, and we also record keyboard events for analysis (not used as model inputs at evaluation time). Captions for CLIGen (General) are generated in three tiers (semantic, regular, detailed) using Llama~3.1~70B to standardize narration. The second mode, CLIGen (Clean), is collected using the open-source \texttt{vhs}\footnote{\url{https://github.com/charmbracelet/vhs}} toolkit (``write terminal GIFs as code''), which lets us author deterministic command scripts that drive Dockerized environments to capture clean, debounced traces (package installs, log filters, interactive REPLs, etc.). For CLIGen (Clean), captions are derived directly from the \texttt{vhs} scripts to provide literal, deterministic supervision. Rendering is fixed by using one monospace font/size, a consistent success/error palette, and a fixed resolution/theme. Each episode records its caption type and font settings for later slicing. We release both splits and run training and analysis per split rather than mixing them.

\noindent\textbf{\ncguiworld{} traces.}
Desktop data comes from screen capture with synchronized pointer and key logs. We collect two random-interaction styles: slow (about 1{,}000 hours) and fast (about 400 hours) to cover cursor motion, window transitions, and idle gaps. The slow style exposes drift after long pauses, while the fast style stresses acceleration and hover timing. Around 110 hours of supervised trajectories from Claude CUA complement this with higher-signal action--response pairs, anchoring the action head without overwhelming exploration. Every sample records the active conditioning mode (\texttt{none}, \texttt{residual}, \texttt{contextual}, \texttt{internal}, or \texttt{external}), letting us slice metrics by the amount of structure given to the latent. Optional accessibility cues (e.g., spoken hints) are logged for analysis but never required by the model; they help diagnose failures such as missed tooltips or misread dialogs.

All GUI capture runs in an Ubuntu~22.04 container with XFCE4 (Arc-Dark, Papirus) on a fixed 1024$\times$768 virtual display at 15~FPS. Xvfb provides the display, x11vnc exposes port~5900, and noVNC exposes 6080; sessions run as user \texttt{computeruse} on \texttt{DISPLAY{:}1}. The desktop pins a small open-source set to launchers (Firefox ESR, GIMP, VLC, VS~Code, Calculator, Terminal, the file manager, and the Mahjongg game), matching the environment shown in our recordings. Screen capture uses \texttt{mss}+\texttt{ffmpeg} with cursor overlays; actions are replayed and logged via \texttt{xdotool}, producing MP4s and CSV/JSON logs under \texttt{/home/computeruse/agent\_recordings}. Synthetic traces use the same stack, generating 30~s, 15~FPS Bezier/random trajectories in parallel after committing a clean base container; real CUA runs use the Anthropic computer-use toolchain with the same display and recording setup.

Unlike joystick-controlled games, GUI actions are high-dimensional and event-driven. Pointer coordinates, scroll deltas, button states, and key presses arrive as separate channels rather than continuous control vectors. Operating systems deliver events in bursts, so the timeline is jagged rather than smooth; pointer updates often appear as step changes after OS-level coalescing, and we log cursor trajectories explicitly instead of inferring them from the video. Keyboard bursts can hit 8--10 characters in a 100~ms window. We keep these discontinuities and the 30--80~ms of interface latency intact, since smoothing them caused the model to hallucinate interpolated motion or underestimate typing speed.

\noindent\textbf{Alignment and preprocessing.}
Frames, text buffers, and event logs share a monotonic clock; we correct residual drift with local cross-correlation between pointer deltas and motion fields. Without this correction the model latches onto stale pointer positions and hallucinates phantom clicks. We normalize resolution and aspect ratio, while privacy filters redact emails, IDs, and faces before persistence. We define a unified action schema across CLI and GUI (character insertions, special keys, pointer moves, clicks, scrolls): in this paper we use it for GUI conditioning and log CLI keystrokes for analysis, and it enables future joint models to share an action vocabulary across interfaces. Traces are windowed into overlapping episodes to cover different interaction densities, and the windows double as curriculum knobs during training. Per-episode identifiers, perceptual hashes, and metadata (application, theme, resolution, caption tier/type, action mode) travel with the data to keep splits and audits deterministic.

\noindent\textbf{Sampling and splits.}
Training draws a balanced mix of short and medium horizons so models stay faithful locally yet stable across multiple steps; overemphasizing either extreme can collapse long-term coherence or blur fine motion. We maintain disjoint train/validation/test splits by application and scenario to avoid interface-specific leakage, and we publish the hash of every split so reviewers can reproduce the exact partitions. Unless otherwise noted, reported metrics are computed on held-out episodes with fixed evaluation code. Weak labels such as OCR transcripts or component boundaries can be attached for ablations, but the core models remain pixel-grounded so improvements reflect the NC rather than auxiliary classifiers.

\noindent\textbf{Quality assurance and privacy.}
Duplicate detection, frozen-screen checks, pointer-motion discordance alerts, and font/rendering audits run continuously. Episodes that desynchronize beyond a small threshold are dropped; keeping bad rollouts hurts training more than losing a little data. All logs pass through sensitive-pattern filters, and only sanitized splits are released.

\noindent\textbf{Release artifacts.}
We release preprocessing scripts, rendering and logging configs, and minimal launchers that reconstruct the data splits used in this paper. Every script accepts a fixed random seed and emits split hashes so others can audit the pipeline end-to-end. In practice: keep every modality on one clock, correct drift early, guard splits from leakage, and prefer simple tooling when quality is equivalent. Violating any of these rules degraded performance more than scaling the model down.

\begin{smallschemetable}{Signals and collection settings.}{tab:data-signals}
  \begin{tabularx}{\linewidth}{L{2.1cm} Y C{2.6cm} Y}
    \toprule
    Interface & Signals & Typical rate & Notes \\
    \midrule
    CLI & frames, text, keystrokes & 15~FPS (resampled) & terminal buffer rendering; cursor/focus \\
    GUI & frames, pointer, keys & 15~FPS & window metadata; optional accessibility cues \\
    \bottomrule
  \end{tabularx}
\end{smallschemetable}
